\chapter{致\hskip\ccwd{}谢}

% 这里用盲审环境包裹致谢，在开启盲审开关时，环境内部的内容不予渲染。

\begin{secretizeEnv}
行文至此，本科毕业论文已然告一段落，而思考还将继续。在完成此论文的过程中，离不开许多人的帮助与支持，谨在此致以诚挚的感谢。

感谢我的导师罗远新教授。从选题到最终成稿，罗老师给予了充分的自由度，允许我在一个不那么"传统"的方向上认真探索，并在关键节点上给出了准确的方向性判断，帮助我少走了许多弯路。在项目推进过程中，面对来自各方的质疑与不同声音，罗老师一次次在价值层面给予的肯定与支持，让我得以更加笃定地走在这条路上。

同时，作为明月科创实验班的创办者，罗老师在新工科教育改革中所展现出的韧性与远见，值得我长期学习。加入明月班，是我大学期间做过的最正确的决定之一。明月班给了我很多机会，但更珍贵的，是让我学会了用一种不同的方式去理解这个世界——理解工程，理解创新，也理解自己。感谢明月班当初选择了我，重塑了我观察世界的视角。

感谢张雷副教授在论文写作过程中给予的专业指导与帮助，让本文在技术深度与规范性上得到了显著提升。同时感谢重庆大学国家卓越工程师学院的各位老师，感谢你们的悉心培养，为我们构建了一片可以自由生长的土壤。

感谢东莞远创模型开发有限公司，为样机（尤其是无人机飞控部分）的开发提供了专业的教程、指导与建议，让我得以在很短的时间内完成无人机控制系统的入门与上手。

感谢临沂华达英飞科技高老板，凭借极为丰富的行业经验，对样机的BOM选型给予了专业的建议与评估，大大提升了方案的可靠性。

感谢来自互联网的老学长 Zhennan LI，以及其开源的 CQUThesis 论文模板——正是因为有人愿意把这些经验整理好留下来，后来者才得以把更多精力放在内容本身上。同时感谢 \LaTeXe{} 计划与 CTeX 开发组，为中文排版提供了完善的解决方案。

感谢我的父母，在整个大学阶段始终给予我无条件的支持——无论是物质上的投入，还是每一次迷茫与低谷时的陪伴与鼓励。你们是我能够放手去做这一切最坚实的底气。

感谢一路同行的朋友们，感谢你们的理解、陪伴与真诚，让这段求学经历有了更多温度与意义。

世界因我们更美好。
\end{secretizeEnv}
